大型语言模型（LLM）在从人类或缺乏长链思维能力（Long CoT）的模型模仿中，往往难以学会有效的长链式思维链推理。为理解这一现象，我们提出：真正有效且可学习的长链思维轨迹在统一视角下呈现出类似分子的稳定结构，由三类交互键构成：深层推理（类共价键）、自我反思（类氢键）与自我探索（类似范德华力）。对蒸馏轨迹的分析显示，这些结构来自长链思维微调本身，而非对关键词的肤浅模仿。我们进一步提出“有效语义异构体（Effective Semantic Isomers）”，并证明只有能促使熵快速收敛的键型才能支撑稳定的长链思维学习，而结构竞争会破坏训练。基于此发现，我们提出 \method，这是一种分布迁移图（distribution-transfer-graph）方法，用来引导合成有效的长链思维结构，可在多个基准上提升性能并增强强化学习过程的稳定性。